\documentclass[11pt, a4paper]{article}

\usepackage[utf8]{inputenc}
\usepackage[T1]{fontenc}
\usepackage[ngerman]{babel}
\usepackage{amsmath}
\usepackage{amssymb}
\usepackage{graphicx}
\usepackage{geometry}
\usepackage{fancyhdr}
\usepackage{enumitem}
\usepackage{array}
\usepackage{eso-pic}

\usepackage{siunitx}
\DeclareSIUnit{\litre}{l}
\DeclareSIUnit{\barPr}{bar}

\sisetup{angle-symbol-degree = ^\circ}
\sisetup{locale = DE, separate-uncertainty, inter-unit-product = {},
range-units = brackets, list-units = single, per-mode=symbol-or-fraction}

\makeatletter
\def\input@path{{./}{../../}}
\makeatother
\graphicspath{{./}{../../}}

\geometry{a4paper, top=2.5cm, bottom=2.5cm, left=2.5cm, right=2.5cm}

\input{defs}

\pagestyle{fancy}
\fancyhf{}
\cfoot{\thepage}

\begin{document}
\noindent
\vspace{9cm}
\begin{center}
    {\Huge \textbf{Lösung}} \\[1.5em]
    {\Large \textbf{Physikalische Grundlagen}} \\[1em]
    {\large 3. Schriftliche Prüfung: Sommersemester 2026} \\[1em]
    {\large Studiengang: Clinical Engineering}
\end{center}

\newpage

\section*{Lösung Beispiel 1 -- Wärmepumpe und Kältemittel}
\begin{enumerate}
    \item \textbf{Temperatur am Ende des Verdampfungsprozesses:}
    Während des Verdampfens (Phasenübergang) bleibt die Temperatur konstant. Daher gilt am Ende des vollständigen Verdampfens noch immer
    \[
        T_1 = \SI{-5}{\degreeCelsius}
    \]

    \item \textbf{Gesamte Verdampfungswärme:}
    Es gilt
    \[
        Q_{\text{Verd}} = m \lambda_{\text{Verd}} = \SI{2.0}{\kilogram} \cdot \SI{210}{\kilo\joule\per\kilogram} = \SI{420}{\kilo\joule}
    \]
    Diese Wärmemenge wird der kalten Umgebung entzogen. Umgerechnet in \si{\kilo\watt\hour} ergibt sich
    \[
        Q_{\text{Verd}} = \frac{\SI{420}{\kilo\joule}}{\SI{3600}{\second\per\hour}} \approx \SI{0.117}{\kilo\watt\hour}
    \]

    \item \textbf{Druck $p_2$ bei isochorer Zustandsänderung:}
    Zuerst werden die Temperaturen in Kelvin umgerechnet:
    \[
        T_1 = \SI{-5}{\degreeCelsius} = \SI{268.15}{\kelvin}
    \]
    \[
        T_2 = \SI{70}{\degreeCelsius} = \SI{343.15}{\kelvin}
    \]
    Für ein ideales Gas bei konstantem Volumen gilt
    \[
        \frac{p}{T} = \const
    \]
    also
    \[
        \frac{p_1}{T_1} = \frac{p_2}{T_2}
    \]
    und damit
    \[
        p_2 = p_1 \frac{T_2}{T_1}
    \]
    Erst jetzt werden die Zahlenwerte eingesetzt:
    \[
        p_2 = \SI{260}{\kilo\pascal} \cdot \frac{\SI{343.15}{\kelvin}}{\SI{268.15}{\kelvin}} \approx \SI{332.7}{\kilo\pascal} \approx \SI{3.33}{\barPr}
    \]

    \item \textbf{Grobe Abschätzung der Verdichterenergie:}
    In diesem vereinfachten isochoren Modell schätzen wir die mindestens nötige Energie über die Änderung der inneren Energie ab:
    \[
        \Delta U = m c_V \Delta T
    \]
    mit
    \[
        \Delta T = T_2 - T_1 = \SI{70}{\degreeCelsius} - \SI{-5}{\degreeCelsius} = \SI{75}{\kelvin}
    \]
    Erst jetzt werden die Zahlenwerte eingesetzt:
    \[
        \Delta U = \SI{2.0}{\kilogram} \cdot \SI{0.85}{\kilo\joule\per(\kilogram\cdot\kelvin)} \cdot \SI{75}{\kelvin} = \SI{127.5}{\kilo\joule}
    \]
    Die Verdichterenergie wird also grob abgeschätzt zu
    \[
        E_{\text{Verdichter}} \approx \Delta U \approx \SI{127.5}{\kilo\joule}
    \]
    \textit{Bemerkung:} Dies ist eine vereinfachte Abschätzung. Ein realer Verdichter arbeitet nicht streng isochor, und die tatsächliche benötigte Energie hängt vom realen Prozessverlauf und weiteren Verlusten ab.
\end{enumerate}

\newpage

\section*{Lösung Beispiel 2 -- Statischer Auftrieb mit Heliumballons}
\begin{enumerate}
    \item \textbf{Nettoauftriebskraft eines einzelnen Ballons:}
    Zuerst wird das Volumen in SI-Einheiten geschrieben:
    \[
        V_{\text{Ballon}} = \SI{4.5}{\litre} = \SI{0.0045}{\metre\cubed}
    \]
    Die Auftriebskraft nach Archimedes lautet
    \[
        F_A = \rho_{\text{Luft}} V_{\text{Ballon}} g
    \]
    \[
        F_A = \SI{1.23}{\kilogram\per\metre\cubed} \cdot \SI{0.0045}{\metre\cubed} \cdot \SI{9.81}{\metre\per\second\squared} \approx \SI{0.0543}{\newton}
    \]
    Die Gewichtskraft des gefüllten Ballons ist
    \[
        F_{G,\text{ges}} = (m_{\text{Hülle}} + \rho_{\text{He}} V_{\text{Ballon}}) g
    \]
    \[
        F_{G,\text{ges}} = (\SI{0.003}{\kilogram} + \SI{0.18}{\kilogram\per\metre\cubed} \cdot \SI{0.0045}{\metre\cubed}) \cdot \SI{9.81}{\metre\per\second\squared} \approx \SI{0.0374}{\newton}
    \]
    Damit ergibt sich die Nettoauftriebskraft:
    \[
        F_{\text{netto}} = F_A - F_{G,\text{ges}} \approx \SI{0.0543}{\newton} - \SI{0.0374}{\newton} = \SI{0.0169}{\newton}
    \]

    \item \textbf{Benötigte Anzahl von Ballons:}
    Die Gewichtskraft der Person ist
    \[
        F_{G,\text{Person}} = mg = \SI{80}{\kilogram} \cdot \SI{9.81}{\metre\per\second\squared} = \SI{784.8}{\newton}
    \]
    Die benötigte Anzahl Ballons ist daher
    \[
        N = \frac{F_{G,\text{Person}}}{F_{\text{netto}}}
    \]
    \[
        N = \frac{\SI{784.8}{\newton}}{\SI{0.0169}{\newton}} \approx 4.64 \cdot 10^4
    \]
    Man benötigt also rund
    \[
        N \approx 46363
    \]
    Ballons.

    \item \textbf{Maximale Steighöhe:}
    Mit zunehmender Höhe nimmt die Dichte der Luft ab. Dadurch wird auch die Auftriebskraft kleiner, da
    \[
        F_A = \rho_{\text{Luft}}(h) V g
    \]
    gilt. Der Aufstieg endet, wenn die Auftriebskraft gerade gleich der Gewichtskraft ist.

    \item \textbf{Vergleich Atmosphäre und Wasser:}
    In der Atmosphäre nimmt die Dichte mit der Höhe deutlich ab. In Wasser bleibt die Dichte mit zunehmender Tiefe dagegen näherungsweise fast konstant und ändert sich nur sehr wenig. Deshalb nimmt der Auftrieb in Luft mit der Höhe wesentlich stärker ab als der Auftrieb im Wasser mit der Tiefe zu- oder abnimmt.

    \item \textbf{Ballon unter Wasser:}
    Beim Aufsteigen nimmt der Umgebungsdruck ab. Dadurch dehnt sich die Luft im Ballon aus, sein Volumen wird größer und damit steigt auch die Auftriebskraft. Der Ballon steigt deshalb immer schneller auf. Hält das Material die Volumenzunahme nicht aus, kann der Ballon platzen.
\end{enumerate}

\newpage

\section*{Lösung Beispiel 3 -- Planet auf Kreisbahn}
\begin{enumerate}
    \item \textbf{Umlaufdauer in Sekunden:}
    \[
        T = \SI{300}{\day} = 300 \cdot 24 \cdot 3600 \si{\second} = \SI{2.592e7}{\second}
    \]

    \item \textbf{Formel für die Bahngeschwindigkeit:}
    Der Bahnumfang einer Kreisbahn ist
    \[
        U = 2\pi r
    \]
    Daher folgt für die Bahngeschwindigkeit
    \[
        v = \frac{U}{T} = \frac{2\pi r}{T}
    \]

    \item \textbf{Bahngeschwindigkeit des Planeten:}
    \[
        v = \frac{2\pi \cdot \SI{2.25e11}{\metre}}{\SI{2.592e7}{\second}} \approx \SI{5.45e4}{\metre\per\second}
    \]
    also
    \[
        v \approx \SI{54.5}{\kilo\metre\per\second}
    \]

    \item \textbf{Formeln für die Zentripetalbeschleunigung:}
    Einerseits gilt
    \[
        a_{\text{Zp}} = \frac{v^2}{r}
    \]
    Setzt man $v = \frac{2\pi r}{T}$ ein, so folgt
    \[
        a_{\text{Zp}} = \frac{1}{r} \left(\frac{2\pi r}{T}\right)^2 = \frac{4\pi^2 r}{T^2}
    \]

    \item \textbf{Numerischer Wert der Zentripetalbeschleunigung:}
    \[
        a_{\text{Zp}} = \frac{v^2}{r} = \frac{(\SI{5.45e4}{\metre\per\second})^2}{\SI{2.25e11}{\metre}} \approx \SI{1.32e-2}{\metre\per\second\squared}
    \]
    also
    \[
        a_{\text{Zp}} \approx \SI{0.0132}{\metre\per\second\squared}
    \]

    \item \textbf{Zentripetalkraft auf den Planeten:}
    \[
        F_{\text{Zp}} = m_{\text{P}} a_{\text{Zp}}
    \]
    \[
        F_{\text{Zp}} = \SI{6.0e24}{\kilogram} \cdot \SI{1.32e-2}{\metre\per\second\squared} \approx \SI{7.93e22}{\newton}
    \]
\end{enumerate}

\newpage

\section*{Lösung Beispiel 4 -- Single-Choice}
\begin{enumerate}
    \item \textbf{Leistung}
    \item \textbf{Das Volumen}
    \item \textbf{Faktor 4}
    \item \textbf{dem Gewicht des verdrängten Mediums}
    \item \textbf{Kräfte treten paarweise als gleich große, entgegengesetzt gerichtete Wechselwirkungen auf.}
    \item \textbf{\SI{3}{\metre\per\second}} \quad denn es gilt $v = r\omega$.
    \item \textbf{Die maximale Haftreibung ist im Allgemeinen größer als die Gleitreibung.}
    \item \textbf{mittlere kinetische Energie der Teilchen}
    \item \textbf{Energieerhaltungssatzes}
    \item \textbf{Arbeit, Einheit: \si{\joule}}
\end{enumerate}

\end{document}